\chapter*{On the Empirical Results}
\addcontentsline{toc}{chapter}{On the Empirical Results}

All numerical claims in this dissertation are intended to be script-derived and
manifest-backed.  The manuscript therefore separates \emph{universal runtime
architecture} from \emph{domain promotion status}.  A row appears as
\emph{reference}, \emph{proof-validated}, \emph{proof-candidate},
\emph{experimental}, or \emph{portability-only} according to the evidence
artifacts that exist for that row, not according to rhetorical ambition.

The battery energy-management row is the deepest witness.  Industrial process
control and medical monitoring are defended proof-validated peers on the
current artifact surface.  Autonomous vehicles remain proof-candidate,
aerospace remains experimental, and navigation remains portability evidence
until the missing replay and promotion surfaces are closed.  Readers should
therefore interpret cross-domain tables as a tiered validation program rather
than as a flat statement of equal maturity across six domains.

Where the manuscript reports confidence intervals, theorem gates, parity
matrices, or release-manifest checks, those objects govern what the surrounding
prose is allowed to say.  Projected numbers, inferred parity, and unsupported
cross-domain closure language are intentionally excluded from the defended
surface.
